% part: intuitionistic-logic
% chapter: propositions-as-types

\documentclass[../../../include/open-logic-chapter]{subfiles}

\begin{document}

\olchapter{int}{pty}{પ્રકારો તરીકે પ્રતિજ્ઞાપનો}

\begin{editorial}
  કરી--હાવર્ડ અનુરૂપતા અંગેના અધ્યાયનો આ
  \emph{અત્યંત પ્રાયોગિક} મુસદ્દો છે.
  વધુ ખુલાસા અને પ્રેરણા જોઈએ, અને તેમાં
  સંભવતઃ ભૂલો તથા ખૂટતી વિગતો છે.
  સામાન્યીકરણની સાબિતીની સમીક્ષા કરીને
  તેને વિસ્તૃત કરવી જોઈએ. ગુણાકાર-પ્રકારનાં
  ઉદાહરણો નથી. ક્રમચય અને સરળીકરણનાં
  રૂપાંતરો આવરી લીધાં નથી. (પ્રકારિત)
  લૅમ્બડા-કલન વિશેની સામગ્રી પણ મળે
  ત્યારે આ વધુ સમજાશે, કારણ કે અહીં
  મૂળભૂત રીતે તે જાણીતું માની લીધું છે.
  અત્યંત સાવચેતીથી વાપરો.
  \sourcecorrection{OLMD-290}{આ મૂળની સંપાદકીય નોંધ યથાવત્ અનુવાદી છે. જોકે સાબિતીથી પદોમાં રૂપાંતરના વિભાગમાં (\olref[pt]{sec}) ગુણાકાર-પ્રકારની અદલાબદલીનું એક ઉદાહરણ છે અને ઘટાડાના વિભાગમાં (\olref[red]{sec}) ક્રમચય-રૂપાંતરનું એક ઉદાહરણ છે. એટલે આ વિષયોનું સંપૂર્ણ આવરણ નથી, પરંતુ એકેય ઉદાહરણ નથી એવો મૂળનો દાવો હાલની અધ્યાય-સામગ્રી સાથે મેળ ખાતો નથી.}
\end{editorial}

\olimport{introduction}
%\olimport{sequent-natural-deduction}
\olimport{proof-terms}
\olimport{proofs-to-terms}
\olimport{terms-to-proofs}
\olimport{types}
\olimport{reduction}
\olimport{type-preservation}
\olimport{normalization}

\OLEndChapterHook

\end{document}
